\section{Introduction}
\label{sec:introduction}

A freight analyst posing the question ``why don't trucking companies use relay driving?'' at an airport gate receives a peculiar answer from nearly every carrier executive: ``the economics don't work.'' The peculiarity is that the economics work unambiguously and demonstrably in relay's favor. A solo driver on the New York--Los Angeles corridor earns approximately \$1,456 per 2,800-mile trip at the ATRI reported average rate of \$0.52 per mile \citep{ATRI_shortage2024}. A relay team of six drivers, each working a 7-hour shift at \$25 per hour, delivers the same load for \$1,050 --- 28\% less. The truck under relay runs at 98\% utilization because it never sits idle during a mandatory rest stop; under the Hours of Service cycle, a solo driver must rest 10 hours per 21-hour on-duty period, leaving the truck idle 48\% of the time \citep{FMCSA_HOS2022}. At \$180,000 per new Class 8 tractor, the utilization premium alone justifies a fleet transformation: a relay-operating carrier needs roughly half the truck capital per unit of freight-ton-miles delivered.

These numbers are not contested in industry. What is contested --- or more precisely, what is never quite articulated --- is whether the coordination cost of relay is worth bearing. The carrier executive at the gate is not wrong that relay is complex. But complexity is not economics. The executive conflates the cost of a missing platform with a structural economic disadvantage. The airport the executive is sitting in resolves an incomparably more complex coordination problem --- gate assignments, crew scheduling, fuel, ground handling, baggage transfer, air traffic control handoffs --- for thousands of flights per day, across hundreds of airlines, none of which owned the infrastructure they use. What the airline industry built between 1950 and 1980 was a platform. Trucking has not built its equivalent.

This paper is about that platform. We call it the relay marketplace: the set of institutional, technological, and regulatory components that together make inter-company relay driving possible for the 70,000+ carriers in the US market \citep{FHWA_freight2023}. The relay marketplace is not a single company or a single technology. It is a four-layer architecture: (1) a standardized truck interface (the Relay Display Unit), (2) an insurance and liability framework that travels with the driver rather than the carrier, (3) a hub slot system that allocates scarce dock capacity across Tier 1, Tier 2, and Tier 3 carriers, and (4) a marketplace platform that matches arriving trucks to available relay drivers and settles the transaction. Each layer is independently necessary; none is sufficient alone.

\subsection*{The Problem in Aggregate}

The scale of the coordination failure is large. The Federal Highway Administration estimates that the US freight system moves approximately 11 billion tons of freight per year, of which 70.5\% moves by truck \citep{FHWA_freight2023}. Long-haul truckload freight --- loads traveling more than 500 miles, the segment most amenable to relay --- accounts for approximately 18\% of total truck ton-miles. By ATRI's estimate, driver wages and benefits are 38\% of all-in trucking cost \citep{ATRI_shortage2024}. A 28\% reduction in driver cost on long-haul loads would reduce total trucking cost on that segment by approximately 10.6\%. Applied to the FHWA estimate of \$320 billion in annual long-haul truckload freight revenue, the relay efficiency gain is a \$33 billion annual reduction in freight cost --- a gain that currently goes unrealized because no relay marketplace exists.

At the same time, the Bureau of Labor Statistics reports approximately 3.5 million truck driver positions in the US, with the long-haul segment facing a structural shortage of 60,000--80,000 drivers as of 2023 \citep{BLS_trucking2023}. The shortage is concentrated in the solo long-haul segment: drivers who must spend weeks away from home, manage their own HOS compliance, and accept the health consequences of the irregular sleep schedule endemic to long-haul solo driving. Regional relay drivers --- who work predictable 7-hour shifts, sleep at home, and never leave their hub region --- do not face the same shortage. The relay marketplace, by design, draws from the regional driver pool rather than the long-haul pool. It does not compete for a scarce resource; it substitutes an abundant one.

\subsection*{The Simulation Evidence}

The ROUTE sla-matrix simulation \citep{ROUTE_C1}, running 5,000 Monte Carlo trips per corridor on the Interstate 2.0 candidate network, provides the quantitative benchmark for relay's contribution to freight reliability. The primary result: relay driving on existing infrastructure reduces NY--LA p95 transit time from 87.2 hours to 58 hours. The full Interstate 2.0 managed-lane program reduces NY--LA p95 to 45.4 hours. The relay-only improvement --- achieved at a capital cost of \$40 million for 8 relay stations on the NY--LA corridor, versus \$121 billion for managed-lane construction \citep{ROUTE_E1} --- captures 90\% of the total SLA improvement at 0.02\% of the infrastructure cost.

The implication is not that managed lanes are unnecessary. The managed-lane program \citep{ROUTE_E2} delivers permanent freight capacity, PTI reduction, platooning enablement, and the physical AV operating environment; these benefits are irreducible to relay alone. The implication is sequencing: relay is the first and cheapest intervention, achievable within the current administration cycle, requiring no new highway construction and no eminent domain. It can be operational in 18--24 months with the right regulatory changes. Managed lanes take 10--15 years. The relay marketplace buys the 15 years of managed-lane construction.

\subsection*{Contributions}

This paper makes four contributions to the platform design and transportation economics literature:

\begin{enumerate}
  \item \textbf{Coordination failure anatomy.} We provide the first systematic analysis of why relay driving has not emerged at scale despite its economic superiority, identifying regulatory gap, first-mover disadvantage, carrier fragmentation, and absent standardization as the four interlocking barriers (Section~\ref{sec:why-relay-fails}).

  \item \textbf{The container model for trucks.} We propose the Relay Display Unit (RDU) cab standard --- a truck standardization layer analogous to the ISO shipping container that McLean \citep{McLean1956} introduced in 1956 --- as the foundation of interoperability across carriers (Section~\ref{sec:container-model}).

  \item \textbf{Dual-mode insurance architecture.} We specify a two-mode insurance framework (Mode 1 hub-operator W-2 employment; Mode 2 independent contractor gig layer) that maps to existing regulatory categories and identifies the specific FMCSA rulemaking required for Mode 2 activation (Section~\ref{sec:insurance}).

  \item \textbf{Hub slot mechanism design.} We specify the three-tier slot allocation system, demonstrate Atlanta hub capacity math (28,740 trucks/day serviced by 750 dock spaces at 20-minute average dwell), and show that the relay marketplace \textit{lowers} effective barriers to entry for small carriers relative to the current solo long-haul model (Section~\ref{sec:hub-slot}).
\end{enumerate}

The paper proceeds as follows. Section~\ref{sec:why-relay-fails} diagnoses the coordination failure in detail. Section~\ref{sec:container-model} introduces the RDU standardization framework. Section~\ref{sec:insurance} specifies the insurance architecture. Section~\ref{sec:hub-slot} develops the slot system. Section~\ref{sec:marketplace} describes the full marketplace platform. Section~\ref{sec:av-transition} examines the autonomous vehicle transition case. Section~\ref{sec:conclusion} states policy recommendations and the sequencing argument.
